
\begin{abstract}
% Recent approaches based on neural networks can learn feature representation automatically from data and achieve state-of-the-art performance on various sequence tagging datasets. However, it is not clear whether the representation learned from one task can be generally useful for another task. In this work, we explore the transfer learning setting where we aim to transfer the knowledge (i.e., model parameters and feature representation) from a source task to a target task so that we can improve the model performance on the target task. We examine the effects of transferring between domains, applications, and languages based on a deep hierarchical recurrent network. We empirically show that our transfer learning approach achieves significant improvement on various low-resource conditions, as well as new state-of-the-art results on several datasets.

Recent papers have shown that neural networks obtain state-of-the-art performance on several different sequence tagging tasks. One appealing property of such systems is their generality, as excellent performance can be achieved with a unified architecture and without task-specific feature engineering.  However, it is unclear if such systems can be used for tasks without large amounts of training data. In this paper we explore the problem of transfer learning for neural sequence taggers, where a source task with plentiful annotations (e.g., POS tagging on Penn Treebank) is used to improve performance on a target task with fewer available annotations (e.g., POS tagging for microblogs). We examine the effects of transfer learning for deep hierarchical recurrent networks across domains, applications, and languages, and show that significant improvement can often be obtained.  These improvements lead to improvements over the current state-of-the-art on several well-studied tasks.\footnote{Code is available at \url{https://github.com/kimiyoung/transfer}}

% NER for Dutch or 
% NER on low-resource languages or 
% -- in particular, recurrent neural networks that include both character-level and word-level embeddings often perform extremely well.  
\end{abstract}
